\subsection{HAUSDORFF GROUP}

Let \(\mathbf{X}\) be a set. We use the notation of § 5, nos. 1 and 2. The free Lie algebra \(\mathrm{L}(\mathbf{X})\) is identified with its canonical image in \(\mathrm{A}(\mathbf{X})\) (§ 3, no. 1, Theorem 1). We denote by \(\hat{\mathrm{L}}(\mathbf{X})\) the closure of \(\mathrm{L}(\mathbf{X})\) in \(\hat{\mathrm{A}}(\mathbf{X})\) , that is the set of elements of \(\hat{\mathrm{A}}(\mathbf{X})\) of the form \(a = \sum_{n\geqslant 1}a_n\) such that \(a_{n}\in \mathrm{L}^{n}(\mathrm{X})\) for all \(n\geqslant 0\) ; this is filtered Lie subalgebra of \(\hat{\mathrm{A}}(\mathbf{X})\) .

THEOREM 1. The restriction of the exponential mapping of \(\hat{\mathbf{A}}(\mathbf{X})\) to \(\hat{\mathbf{L}}(\mathbf{X})\) is a bijection of \(\hat{\mathbf{L}}(\mathbf{X})\) onto a closed subgroup of the Magnus group \(\Gamma(\mathbf{X})\) .

We write \(\mathrm{A}(\mathrm{X}) = \mathrm{A}\) , \(\mathrm{A}^{\mathrm{n}}(\mathrm{X}) = \mathrm{A}^{\mathrm{n}}\) , \(\hat{\mathrm{A}}(\mathrm{X}) = \hat{\mathrm{A}}\) , \(\mathrm{L}^{\mathrm{n}}(\mathrm{X}) = \mathrm{L}^{\mathrm{n}}\) , \(\hat{\mathrm{L}}(\mathrm{X}) = \hat{\mathrm{L}}\) , \(\Gamma(\mathrm{X}) = \Gamma\) . Let \(\mathrm{B}\) be the algebra \(\mathrm{A} \otimes \mathrm{A}\) with the graduation of type \(\mathbf{N}\) defined by \(\mathrm{B}^{\mathrm{n}} = \sum_{i+j=n} \mathrm{A}^{i} \otimes \mathrm{A}^{j}\) . Let \(\hat{\mathrm{B}} = \prod_{n > 0} \mathrm{B}^{n}\) be the associated complete filtered algebra (Commutative Algebra, Chapter III, §2, no. 12, Example 1). The co-product \(c: \mathrm{A} \to \mathrm{A} \otimes \mathrm{A}\) defined in §3, no. 1, Corollary 1 to Theorem 1 is graded of degree 0 and hence extends by continuity to a homomorphism \(\hat c: \hat{\mathrm{A}} \to \hat{\mathrm{B}}\) given by

\[
\hat {c} \left(\sum_ {n \geqslant 0} a _ {n}\right) = \sum_ {n \geqslant 0} c (a _ {n}) \quad \text { for } a _ {n} \in \mathrm{A} ^ {n}.
\]

We also define continuous homomorphisms \(\delta'\) and \(\delta''\) of \(\hat{\mathbf{A}}\) into \(\hat{\mathbf{B}}\) by

\[
\delta^ {\prime} \left(\sum_ {n \geqslant 0} a _ {n}\right) = \sum_ {n \geqslant 0} (a _ {n} \otimes 1), \quad \delta^ {\prime \prime} \left(\sum_ {n \geqslant 0} a _ {n}\right) = \sum_ {n \geqslant 0} (1 \otimes a _ {n}) \quad \text { for } a _ {n} \in \mathrm{A} ^ {n}.
\]

By Corollary 2 to Theorem 1 of §3, no. 1, \(\mathbf{L}^n\) is the set of \(a_{n}\in \mathbf{A}^{n}\) such that \(c(a_{n}) = a_{n}\otimes 1 + 1\otimes a_{n}\) . It follows that \(\hat{\mathbf{L}}\) is the set of \(a\in \hat{\mathbf{A}}\) such that

\[
\hat {c} (a) = \delta^ {\prime} (a) + \delta^ {\prime \prime} (a).\tag{3}
\]

Let \(\Delta\) be the set of \(b\in \hat{\mathbf{A}}\) of constant term equal to 1 and satisfying the relation

\[
\hat {c} (b) = \delta^ {\prime} (b). \delta^ {\prime \prime} (b),\tag{4}
\]

in other words, the set of \(b = \sum_{n\geqslant 0}b_n\) such that \(b_{n}\in \mathbf{A}^{n}\) for all \(n\geqslant 0\) , \(b_0 = 1\) and \(c(b_{n}) = \sum_{t + j = n}b_{t}\otimes b_{j}\) for \(n\geqslant 0\) . The latter characterization shows that \(\Delta\) is a closed subset of \(\Gamma\) ; as \(\hat c,\delta^{\prime}\) and \(\delta ''\) are ring homomorphisms and every element of \(\delta'(\hat{\mathbf A})\) commutes with every element of \(\delta''(\hat{\mathbf A})\) , the restrictions to \(\Gamma\) of the mappings \(c\) and \(\delta ' \delta''\) are group homomorphisms and \(\Delta\) is a subgroup of \(\Gamma\) .

By Proposition 1 of no. 1, the exponential mapping of \(\hat{\mathbf{A}}\) is a bijection of the set \(\hat{\mathbf{A}}^{+}\) of elements of \(\hat{\mathbf{A}}\) with no constant term onto \(\Gamma\) . Let \(a \in \hat{\mathbf{A}}^{+}\) and \(b = \exp a\) . As \(\hat c\) is a continuous ring homomorphism,

\[
\hat {c} (b) = \hat {c} \left(\sum_ {n \geqslant 0} a ^ {n} / n!\right) = \sum_ {n \geqslant 0} \hat {c} (a) ^ {n} / n! = \exp \hat {c} (a).
\]

The relations

\[
\delta^ {\prime} (b) = \exp \delta^ {\prime} (a), \quad \delta^ {\prime \prime} (b) = \exp \delta^ {\prime \prime} (a)
\]

are proved similarly and, as \(\delta'(a)\) commutes with \(\delta''(a)\) (no. 1, Remark 1),

\[
\delta^ {\prime} (b) \delta^ {\prime \prime} (b) = \exp (\delta^ {\prime} (a) + \delta^ {\prime \prime} (a)).
\]

Therefore a satisfies (3) if and only if b satisfies (4), which proves the theorem. Remark. The above proof shows that \(\exp(\hat{\mathbf{L}})\) is the subgroup \(\Delta\) of \(\Gamma\) consisting of the b satisfying (4).

Hence the group law of \(\Delta\) can be transported by the exponential mapping to \(\hat{\mathbf{L}}\) . In other words, \(\hat{\mathbf{L}}\) is a complete topological group with the law of composition \((a, b) \mapsto a \uplus b\) given by

\[
a \uplus b = \log (\exp a. \exp b).
\]

The topological group thus obtained is called the Hausdorff group (derived from X relative to K).

Let \(g\) be the homomorphism of the free group \(\mathbf{F} = \mathrm{F}(\mathbf{X})\) into \(\Gamma\) such that \(g(x) = \exp x\) for \(x \in \mathbf{X}\) . As \(\exp x - 1 - x = \sum_{n \geqslant 2} x^n / n!\) is of order \(\geqslant 2\) , \(g\) is injective by Theorem 1 of §5, no. 3. Therefore the mapping \(\log \circ g\) is an injective homomorphism of \(F\) into the Hausdorff group which extends the canonical injection \(X \to L\) .

For every integer \(m \geqslant 1\) we denote by \(\hat{\mathbf{L}}_m\) the set of elements of order \(\geqslant m\) in \(\hat{\mathbf{L}}\) and by \(\Gamma_m\) the set of \(u \in \Gamma\) such that \(u - 1\) is of order \(\geqslant m\) . Then \(\hat{\mathbf{L}}_m = \exp^{-1}(\Gamma_m)\) by Remark 2 of no. 1; as \((\Gamma_m)_{m \geqslant 1}\) is an integral central filtration on \(\Gamma\) (§4, no. 5 Proposition 2), \((\hat{\mathbf{L}}_m)_{m \geqslant 1}\) is an integral central filtration on the group \(\hat{\mathbf{L}}\) .

